\chapter{On-Chain Connectivity}\label{ch:nw:on-chain-connectivity}

A transaction sent to a public node of a blockchain reaches the rest of the network by gossip, hop by hop, and is visible to everyone on the
way; one sent to a \omterm{def:nw:on-chain-connectivity:engine}{block engine} in the region of the next block producer skips the gossip altogether. On a chain whose leaders change every
few hundred milliseconds, and whose software gives each leader four consecutive slots, which region to send to is a scheduling question: the
answer changes with the \omterm{def:nw:on-chain-connectivity:leader}{leader schedule}, and the schedule is published in advance.

Book~3 described the actors of on-chain trading (nodes, validators, builders, relays, searchers). This chapter is about their network: how
nodes peer and gossip, what propagation networks and private relays change, where the \omterm{def:nw:on-chain-connectivity:engine}{block engines} are, how a \omterm{def:nw:on-chain-connectivity:leader}{leader schedule} turns into a
choice of endpoint, and what \omterm{def:nw:on-chain-connectivity:engine}{hosted node providers} sell. The documentation is cited and dated; the propagation is a labelled simulation on
the book's map.

\section{Running and peering nodes}

\begin{definition}[Gossip protocol]\label{def:nw:on-chain-connectivity:gossip}
A \emph{gossip protocol}\index{gossip protocol} spreads a message over a peer-to-peer network by having each node that receives it for the first
time forward it (or announce it) to some or all of its peers, so that it reaches every node in a number of hops that grows with the logarithm of
the network's size.
\end{definition}

A node joins a chain's network by connecting to peers found through the chain's discovery protocol; in this chapter's model they are chosen at
random, with no regard to geography. Each hop costs the link's delay plus the node's validation of the message before it forwards it. The consequence was measured early:
Decker and Wattenhofer found in 2013 that a Bitcoin block reached the median node after 6.5 seconds and that after 40 seconds 5\% of nodes still
had not seen it. Smaller messages and faster validation shorten each hop, but not the structure: a transaction still reaches the far side of
the world through a chain of hops between peers it did not choose.

\begin{proposition}[Gossip time]\label{prop:nw:on-chain-connectivity:gossip}
On a network whose nodes each forward to all their peers on first receipt, a message reaches every node at the length of its shortest path
from the source, where a link's length is its delay plus the receiving node's processing; forwarding to a subset and announcing to the rest, who
pull, lengthens the announced links by a round trip and can only delay arrival.
\end{proposition}

\begin{proof}
Flooding explores every path in parallel, so each node receives the message first along its shortest path (Dijkstra's argument). Replacing some
pushes by announcements adds a round trip to those links' lengths, and shortest paths cannot become shorter when lengths grow.
\end{proof}

\omcode{code/firm/chainnet/firm_chainnet.py}{54}{71}{Gossip as shortest paths: flooding, and square-root fan-out with announcements and pulls.}\label{lst:nw:on-chain-connectivity:spread}

The model places 600 nodes evenly in the eight cities of \cref{dat:nw:on-chain-connectivity:engines}, joins each to eight random peers, and
charges each link its fibre floor times 1.5 plus \qty{5}{\milli\second} of processing (assumptions). A transaction from a node in Tokyo reaches
half the nodes after \qty{85.0}{\milli\second} and nine-tenths after \qty{99.5}{\milli\second} by flooding; with square-root fan-out, which
saves bandwidth, after 113.9 and \qty{149.0}{\milli\second}. In the model the random graph offers paths close to the direct legs,
so what gossip adds is its extra hops: three to eight of them, 15 to \qty{40}{\milli\second} of processing, against the engine's one
millisecond (an assumption). The tails differ most: the last node has it after \qty{158.5}{\milli\second} by flooding and
\qty{326.5}{\milli\second} with square-root fan-out.

\begin{omfigure}
\begin{tikzpicture}
\begin{axis}[width=11cm,height=5.6cm,xmin=0,xmax=180,ymin=0,ymax=1.02,xlabel={time since the transaction left Tokyo (ms)},
  ylabel={share of nodes reached},tick label style={font=\footnotesize},label style={font=\small},grid=major,grid style={gray!25},
  legend pos=north west,legend style={font=\footnotesize},table/col sep=comma]
\addplot[omDef,thick,const plot] table[x=ms,y=flood]{figdata/networks/21-on-chain-connectivity/propagation.csv};
\addplot[omThm,thick,dashed,const plot] table[x=ms,y=sqrt]{figdata/networks/21-on-chain-connectivity/propagation.csv};
\addplot[gray,densely dotted] coordinates {(0,0.5) (180,0.5)};
\addplot[gray,densely dotted] coordinates {(0,0.9) (180,0.9)};
\legend{flooding,square-root fan-out}
\end{axis}
\end{tikzpicture}

\omcaption{Simulation: the share of 600 nodes in eight cities, eight random peers each, that have a transaction sent from a node in Tokyo, against
time. Flooding reaches half at \qty{85.0}{\milli\second} and nine-tenths at \qty{99.5}{\milli\second}; square-root fan-out at 113.9 and
\qty{149.0}{\milli\second}. Data: \texttt{fig\_chain.py}, \texttt{firm\_chainnet.spread()}.}\label{fig:nw:on-chain-connectivity:propagation}
\end{omfigure}

\begin{omfigure}
\begin{tikzpicture}
\begin{axis}[xbar,width=11cm,height=6.2cm,bar width=6pt,xmin=0,xmax=110,ytick=data,
  yticklabels from table={figdata/networks/21-on-chain-connectivity/leaders.csv}{city},y dir=reverse,
  xlabel={arrival at a leader in the city (ms), from Tokyo},tick label style={font=\footnotesize},label style={font=\small},
  grid=major,grid style={gray!25},legend style={font=\footnotesize,at={(0.5,-0.3)},anchor=north,legend columns=2},
  table/col sep=comma]
\addplot[fill=omThm!45,draw=omThm] table[y expr=\coordindex,x=gossip]{figdata/networks/21-on-chain-connectivity/leaders.csv};
\addplot[fill=omDef!55,draw=omDef] table[y expr=\coordindex,x=direct]{figdata/networks/21-on-chain-connectivity/leaders.csv};
\legend{through gossip,through the nearest block engine}
\end{axis}
\end{tikzpicture}

\omcaption{Simulation: when a transaction sent from Tokyo reaches the next leader, by the leader's city: through the gossip graph from a public node
in Tokyo (median over the city's nodes), and directly through the published \omterm{def:nw:on-chain-connectivity:engine}{block engine} nearest the leader. Gossip adds 14 to
\qty{39}{\milli\second}. Data: \texttt{fig\_chain.py}, \texttt{nw\_chain.to\_leader()}.}\label{fig:nw:on-chain-connectivity:leaders}
\end{omfigure}

\section{Transaction propagation networks and private relays}

\begin{definition}[Transaction propagation network]\label{def:nw:on-chain-connectivity:tpn}
A \emph{transaction propagation network}\index{transaction propagation network} is a privately operated network of relays, peered with many nodes
and validators of a chain, that carries transactions and blocks between them faster and with less variance than the chain's own gossip over the
public internet.
\end{definition}

Two services change the path of a transaction. A propagation network, such as bloXroute's Blockchain Distribution Network, replaces random public
hops with its own relays, which its documentation describes as reducing ``latency and variance'' in propagation. A private relay changes who sees
the transaction: Flashbots Protect, on Ethereum, sends it to a private mempool where it is ``hidden from frontrunning and sandwich bots'' until a
builder includes it. The first is a latency product; the second a visibility product; a searcher may need both (Book~3's private order flow).

\section{Block-builder and block-engine endpoints by region}

\begin{definition}[Block engine, hosted node provider]\label{def:nw:on-chain-connectivity:engine}
A \emph{block engine}\index{block engine} is a service, run close to a chain's validators, that accepts transactions and bundles from traders and
forwards them directly to the current or next block producer, often through an auction for inclusion. A \emph{hosted node provider}\index{hosted node
provider} runs a chain's nodes and sells access to their RPC endpoints, with request limits per plan, so that a firm need not run its own.
\end{definition}

\begin{dated}{2026-09}{Block engines and relays, from their documentation}\label{dat:nw:on-chain-connectivity:engines}
\begin{itemize}
\item Jito (Solana): mainnet \omterm{def:nw:on-chain-connectivity:engine}{block engines} in Amsterdam, Dublin, Frankfurt, London, New York, Salt Lake City, Singapore and Tokyo, each with its own
  endpoint; its \texttt{sendTransaction} forwards a transaction directly to the validator; bundles need a tip of at least 1\,000 lamports.
\item Flashbots Protect (Ethereum): a private mempool hidden from frontrunning and sandwich bots, with refunds.
\item bloXroute: a Blockchain Distribution Network of relays for Solana, BNB Chain, Base, Ethereum, Hyperliquid and other chains.
\end{itemize}
\end{dated}

\begin{omfigure}
\begin{tikzpicture}[font=\small,node distance=6mm,
  nd/.style={circle,draw=omThm,fill=omThm!15,minimum size=5mm,inner sep=0pt},
  box/.style={draw=omDef,fill=omDef!12,rounded corners,align=center,minimum height=8mm,inner sep=3pt}]
\node[box] (firm) at (0,0) {firm\\(Tokyo)};
\node[nd] (a) at (2.2,1.3) {}; \node[nd] (b) at (3.8,1.9) {}; \node[nd] (c) at (5.4,1.2) {};
\node[nd] (d) at (7.0,1.8) {}; \node[nd] (e) at (4.6,0.3) {};
\node[box] (lead) at (9.4,0) {next leader\\(Frankfurt)};
\node[box] (eng) at (7.4,-1.3) {block engine\\(Frankfurt)};
\draw[omThm,->] (firm) -- (a); \draw[omThm,->] (a) -- (b); \draw[omThm,->] (b) -- (c); \draw[omThm,->] (c) -- (d);
\draw[omThm,->] (d) -- (lead); \draw[omThm!50,->] (a) -- (e); \draw[omThm!50,->] (e) -- (c);
\draw[omDef,->,thick] (firm) |- (eng); \draw[omDef,->,thick] (eng) -- (lead);
\node[text=omThm,font=\footnotesize] at (4.6,2.5) {public gossip: every hop sees it and adds its processing};
\node[text=omDef,font=\footnotesize] at (2.9,-1.62) {firm's route to the engine; one short last hop};
\end{tikzpicture}

\omcaption{Two ways to the next leader: hand the transaction to a public node and let gossip carry it hop by hop, visible to every node; or send it
over the firm's own route to the \omterm{def:nw:on-chain-connectivity:engine}{block engine} nearest the leader, which forwards it directly. Schematic.}\label{fig:nw:on-chain-connectivity:paths}
\end{omfigure}

A \omterm{def:nw:on-chain-connectivity:engine}{block engine} is a \omterm{def:nw:the-european-map:pop}{point of presence} (chapter~11) for the chain's producers. Sending to the engine nearest the producer puts the long-haul part of
the path on the firm's own route and the last hop on the engine's short one, instead of on a dozen random public hops.

\section{Validator proximity and the leader schedule}

\begin{definition}[Leader schedule]\label{def:nw:on-chain-connectivity:leader}
A \emph{leader schedule}\index{leader schedule} is a chain's published assignment of future block-production slots to validators, computed in
advance for an epoch, which lets anyone know which validator will produce each coming block and, if its location is known, where to send
transactions.
\end{definition}

\begin{dated}{2026-09}{Slots and leaders, from the chains' specifications and code}\label{dat:nw:on-chain-connectivity:chains}
Ethereum's mainnet configuration sets 12-second slots. Solana computes its \omterm{def:nw:on-chain-connectivity:leader}{leader schedule} an epoch in advance; the Solana SDK's default target
slot time is \qty{300}{\milli\second} (it defines 400 and shorter as well) and the effective target is changed dynamically, so clients must not assume
the default; the Agave validator gives each leader four consecutive slots.
\end{dated}

With 12-second slots, geography matters for Ethereum only at the edges of a slot. With slots of a few hundred milliseconds and four in a row per
leader, it matters for every transaction on Solana: the next leader's location is known in advance, and a transaction that reaches it after its
four slots waits for the next leader, possibly on another continent.

\omcode{code/networks/21-on-chain-connectivity/python/nw_chain.py}{36}{44}{Arrival at the next leader: through the gossip graph, and directly through
the published engine nearest the leader.}\label{lst:nw:on-chain-connectivity:leader}

\begin{omfigure}
\begin{tikzpicture}
\begin{axis}[width=11cm,height=5.6cm,xmin=0,xmax=150,ymin=0,ymax=1.05,xlabel={time left in the leader's window when the transaction is sent (ms)},
  ylabel={share of leaders reached},tick label style={font=\footnotesize},label style={font=\small},grid=major,grid style={gray!25},
  legend pos=north west,legend style={font=\footnotesize},table/col sep=comma]
\addplot[omThm,thick,const plot] table[x=deadline,y=gossip]{figdata/networks/21-on-chain-connectivity/timely.csv};
\addplot[omDef,thick,const plot] table[x=deadline,y=direct]{figdata/networks/21-on-chain-connectivity/timely.csv};
\legend{through gossip,leader-aware engine}
\end{axis}
\end{tikzpicture}

\omcaption{Simulation: the share of next leaders, placed in turn in each of the eight engine cities, reached in time from Tokyo, against the time left
in the leader's window. Sending to the engine nearest the leader shifts the curve left by 14 to \qty{39}{\milli\second}. Data: \texttt{fig\_chain.py},
\texttt{nw\_chain.timely\_rates()}.}\label{fig:nw:on-chain-connectivity:timely}
\end{omfigure}

Averaged over the eight leader cities and over a time left uniform between zero and \qty{150}{\milli\second}, the leader-aware chooser lands 61\% of
transactions in time against 48\% through gossip; over \qty{300}{\milli\second}, 80.5\% against 74.0\%. For a leader in Frankfurt, the
chapter's problem, the direct path takes \qty{69.5}{\milli\second} and gossip \qty{83.5}{\milli\second}.

\section{Hosted node providers}

Running a chain's nodes well is work: storage, bandwidth, upgrades and peering. Hosted
providers sell their nodes' endpoints instead, by plan and request limit. For trading, the question is the one of chapters~16 to~18: where the
provider's nodes are, whether they are the ones the firm's transactions should enter, and how their limits are governed (Book~3's rate limits). A
provider's endpoint is a convenient place to read the chain and a poor place to race on it.

\begin{method}[Getting a transaction to the next block]\label{met:nw:on-chain-connectivity:send}
\begin{enumerate}
\item Read the chain's slot time and \omterm{def:nw:on-chain-connectivity:leader}{leader schedule}; map the coming leaders to regions where their location is known or can be measured
  (chapter~17's triangulation).
\item Keep connections open to the \omterm{def:nw:on-chain-connectivity:engine}{block engines} or builders in every region the leaders use; choose per transaction the one nearest the next
  leader with enough time left.
\item Use a propagation network or private relay where visibility or variance matters more than the last milliseconds.
\item Run your own nodes where you read state critical to trading; use hosted endpoints for the rest, within their limits.
\item Measure landing rates per region and per leader; the schedule and the engines' locations change.
\end{enumerate}
\end{method}

\section{Tutorial: gossip and the next leader}\label{tut:nw:on-chain-connectivity:tut}

\begin{tutorial}
\textbf{Goal.} Simulate gossip on a peer graph over the book's map, and compare it with sending to the \omterm{def:nw:on-chain-connectivity:engine}{block engine} nearest the next leader.
\textbf{End state:} \cref{fig:nw:on-chain-connectivity:propagation,fig:nw:on-chain-connectivity:leaders,fig:nw:on-chain-connectivity:timely}.
\begin{tutsteps}
\item \textbf{Engines.} \texttt{firm\_chainnet.load\_engines} reads the dated regional endpoints.
\item \textbf{Graph.} \texttt{Graph} places 600 nodes in the engines' cities with eight peers each; \texttt{spread} (\cref{lst:nw:on-chain-connectivity:spread})
  gives each node's arrival time.
\item \textbf{Leaders.} \texttt{nw\_chain.to\_leader} (\cref{lst:nw:on-chain-connectivity:leader}) compares gossip with the nearest engine;
  \texttt{timely\_rates} and \texttt{mean\_timely} turn it into landing shares.
\end{tutsteps}
\textbf{What to change next.} Make the graph geography-aware (peers chosen among the nearest nodes) and see what gossip gains; add the engines'
own queueing and tips.
\end{tutorial}

\section{Build: the chain network model}\label{bld:nw:on-chain-connectivity:chainnet}

\begin{build}
\bfield{Purpose} How transactions travel on a chain's network and where to send them: the on-chain rows of chapter~29's plan.
\bfield{Interface} \texttt{firm\_chainnet}: \texttt{Graph}, \texttt{spread}, \texttt{reach\_time}, \texttt{Engine}, \texttt{load\_engines},
\texttt{one\_way\_ms}, \texttt{choose\_engine}, \texttt{timely}.
\bfield{Rules} Endpoints are dated rows with sources; link delays are the fibre floor times a stated \omterm{def:nw:the-north-american-map:factor}{route factor} plus stated processing; every result
is labelled a simulation.
\bfield{Acceptance tests} \texttt{code/firm/chainnet/tests/}: the engines' data, flooding on a line graph by hand, square-root fan-out never faster
than flooding, and the chooser and deadline check.
\bfield{Stretch} Read a live \omterm{def:nw:on-chain-connectivity:leader}{leader schedule} and validators' measured locations; model the engines' auctions.
\end{build}

\begin{omsources}
\item C.~Decker and R.~Wattenhofer, ``Information propagation in the Bitcoin network'', IEEE P2P 2013.
\item Ethereum consensus specifications (mainnet configuration); Anza documentation (leader rotation), Solana SDK and Agave source.
\item Jito Labs, low-latency transaction send; Flashbots Protect; bloXroute documentation.
\end{omsources}

\section{Exercises}

\begin{exercise}[$\star$]\label{exo:nw:on-chain-connectivity:1}
How many milliseconds does a Solana leader hold the chain if it has four consecutive slots at the SDK's default target?
\end{exercise}

\begin{exercise}[$\star$]\label{exo:nw:on-chain-connectivity:2}
What is the difference between a \omterm{def:nw:on-chain-connectivity:tpn}{transaction propagation network} and a private relay?
\end{exercise}

\begin{exercise}[$\star$]\label{exo:nw:on-chain-connectivity:3}
Why does gossip reach a node in the sender's own city later than a direct message would?
\end{exercise}

\begin{exercise}[$\star\star$]\label{exo:nw:on-chain-connectivity:4}
Using \cref{prop:nw:on-chain-connectivity:gossip}, why is square-root fan-out never faster than flooding, and why use it?
\end{exercise}

\begin{exercise}[$\star\star$]\label{exo:nw:on-chain-connectivity:5}
The next leader is in Singapore and your machine in Tokyo. Which engine do you use, and what does the model say you gain?
\end{exercise}

\begin{exercise}[$\star\star$]\label{exo:nw:on-chain-connectivity:6}
Why does the choice of region matter little on a chain with 12-second slots?
\end{exercise}

\begin{exercise}[$\star\star\star$]\label{exo:nw:on-chain-connectivity:7}
\emph{Coding.} With \texttt{firm.chainnet}, how much faster does flooding reach nine-tenths of the nodes if each node has 16 peers instead of 8?
\end{exercise}

\begin{exercise}[$\star\star\star$]\label{exo:nw:on-chain-connectivity:8}
\emph{Find the flaw.} ``We send every transaction to the \omterm{def:nw:on-chain-connectivity:engine}{block engine} nearest our servers, so we always use the fastest path.''
\end{exercise}

\section{Problem: The Leader Is in Frankfurt}

\begin{problem}[{Weekend problem --- gossip against the leader-aware engine}]\label{pb:nw:on-chain-connectivity:1}
A firm in Tokyo sends transactions on a chain whose next leader is in Frankfurt. It can hand them to a public node in Tokyo and let gossip carry them,
or send them to the \omterm{def:nw:on-chain-connectivity:engine}{block engine} nearest the leader. Use the chapter's network model.

\medskip\noindent\textbf{Part I --- Gossip.}
\begin{enumerate}
\item How long does flooding take to reach half and nine-tenths of the nodes from Tokyo?
\item And square-root fan-out?
\item When does the transaction reach a leader in Frankfurt through gossip?
\item Why is that longer than the direct path?
\end{enumerate}

\medskip\noindent\textbf{Part II --- The engine.}
\begin{enumerate}[resume]
\item Which engine is nearest the leader, and when does the transaction reach the leader through it?
\item What is the one-way fibre floor Tokyo--Frankfurt, and what \omterm{def:nw:the-north-american-map:factor}{route factor} does the model assume?
\item How much of the direct path is the engine's own time?
\item What does the engine see that gossip would have shown to everyone?
\end{enumerate}

\medskip\noindent\textbf{Part III --- Landing rates.}
\begin{enumerate}[resume]
\item With the time left uniform between zero and \qty{150}{\milli\second}, what share of leaders over the eight cities does each method reach?
\item And with up to \qty{300}{\milli\second}?
\item Why does the gain shrink as the window grows?
\item How would four consecutive slots per leader change the firm's planning?
\end{enumerate}

\medskip\noindent\textbf{Part IV --- The verdict.}
\begin{enumerate}[resume]
\item State the \emph{named result}: the propagation time to half and nine-tenths of the network, and the gain in timely arrival from sending to the
  \omterm{def:nw:on-chain-connectivity:engine}{block engine} nearest the next leader against a fixed region.
\item Which inputs are documentation and which assumptions?
\item What would a geography-aware peer graph change?
\item Where should the firm's machines be for this chain?
\item What does a propagation network add on top of the engine?
\item How should the firm check that its model matches reality?
\item Where does this go in chapter~29's plan?
\item In one sentence: why is the region a scheduling question?
\end{enumerate}
\end{problem}

\section{Interview questions}

\begin{interviewq}[$\star$ \iqroles{developer}]\label{iq:nw:on-chain-connectivity:1}
How does a transaction propagate on a blockchain's peer-to-peer network, and why is it slow?
\end{interviewq}

\begin{interviewq}[$\star\star$ \iqroles{developer}]\label{iq:nw:on-chain-connectivity:2}
What is a \omterm{def:nw:on-chain-connectivity:engine}{block engine}, and why are there several around the world?
\end{interviewq}

\begin{interviewq}[$\star\star$ \iqroles{developer, researcher}]\label{iq:nw:on-chain-connectivity:3}
How would you use a chain's \omterm{def:nw:on-chain-connectivity:leader}{leader schedule} to route transactions?
\end{interviewq}

\begin{interviewq}[$\star\star$ \iqroles{developer}]\label{iq:nw:on-chain-connectivity:4}
When would you run your own node rather than use a hosted provider?
\end{interviewq}

\begin{interviewq}[$\star\star$ \iqroles{trader}]\label{iq:nw:on-chain-connectivity:5}
Why might a trader send a transaction privately rather than to the public mempool?
\end{interviewq}

\begin{interviewq}[$\star\star\star$ \iqroles{developer, researcher}]\label{iq:nw:on-chain-connectivity:6}
Design the transaction-sending layer of a trading firm active on a fast chain with leaders on three continents.
\end{interviewq}
